\pdfbookmark[2]{Zu IX.}{nexus-remarks-ix}
\begin{center}\large Zu IX.\end{center}

Es sei $p$ eine positive ungerade durch kein Quadrat theilbare Zahl, und
\[
S_r=\Sigma\left(\frac{s_r}p\right)
\]
wo $s_r$ alle relativen Primzahlen zu $p$ durchlaufen muss, welche zwischen $(r-1)\frac p8$ und $r\frac p8$ liegen; bezeichnet man die Anzahlen der nicht eigentlich äquivalenten formae proprie primitivae für die Determinanten $-p$ und $-2p$ resp. mit $C_1$ und $C_2$, so ist (vergl. \textsc{Dirichlet} Recherches sur diverses applications etc. §. 11 in \textsc{Crelle}'s Journal XXI)
\[
C_1=2(S_1+S_2),\quad C_2=2(S_1-S_4)
\]
oder
\[
C_1=S_1+S_2+S_3+S_4,\quad C_2=2(S_2+S_3)
\]
je nachdem $p\equiv1$ oder $\equiv3\pmod4$ ist. Bedenkt man ferner, dass die Zahlen $s_1$ und $s_2$ im Complex mit den Zahlen $2s_1$ und $p-2s_4$, und ebenso die Zahlen $s_3$ und $s_4$ im Complex mit den Zahlen $2s_2$ und $p-2s_3$ übereinstimmen, und dass im Falle $p\equiv1\pmod4$ die Summe $S_1+S_2+S_3+S_4=0$ ist, so ergeben sich in beiden Fällen noch zwei neue Relationen zwischen den vier Summen $S_1,S_2,S_3,S_4$, so dass jede derselben durch $C_1$ und $C_2$ ausgedrückt werden kann. Man erhält auf diese Weise, wenn $p\equiv1\pmod4$ ist,
\[
\begin{aligned}
S_1=S_8&=\tfrac14\left(\frac2p\right)C_1+\tfrac14C_2\\
S_2=S_7&=\tfrac14\left(2-\left(\frac2p\right)\right)C_1-\tfrac14C_2\\
S_3=S_6&=-\tfrac14\left(2+\left(\frac2p\right)\right)C_1+\tfrac14C_2\\
S_4=S_5&=\tfrac14\left(\frac2p\right)C_1-\tfrac14C_2
\end{aligned}
\]
\SourceLeaf{312}
und, wenn $p\equiv3\pmod4$ ist
\[
\begin{aligned}
S_1=-S_8&=\tfrac14\left(3+\left(\frac2p\right)\right)C_1-\tfrac14C_2\\
S_2=-S_7&=-\tfrac14\left(1-\left(\frac2p\right)\right)C_1+\tfrac14C_2\\
S_3=-S_6&=\tfrac14\left(1-\left(\frac2p\right)\right)C_1+\tfrac14C_2\\
S_4=-S_5&=\tfrac14\left(1-\left(\frac2p\right)\right)C_1-\tfrac14C_2
\end{aligned}
\]

Ist $p$ eine Primzahl, so findet man hieraus unmittelbar die im Text angegebenen Formeln für die Anzahl der quadratischen Reste, welche in den einzelnen Octanten enthalten sind.
\BookRule{35mm}
